\section{Notation and conventions}\label{sec:conventions}

\subsection{Notation}

Let $G=(V,E)$ be a graph, with vertex set $V$ and edge set $E$ (the exact
definition is \cref{def:graph}). We use the following letters throughout.

\begin{itemize}[leftmargin=1.5em]
  \item $V=\{n_1,\dots,n_N\}$ is the set of vertices and $N=|V|$ is their
  number. The board names the vertices $n_1,\dots,n_N$ and writes $N$
  \cite{whiteboard1}. The textbook writes $v_1,\dots,v_n$ \cite[p.~11]{petrosyan2015graph}; we do not use
  that form again.
  \item $e_{ij}$ is the edge between $n_i$ and $n_j$. In a simple graph
  $e_{ij}$ and $e_{ji}$ are the same edge.
  \item $A=(A_{ij})$ is the adjacency matrix, an $N\times N$ matrix of zeros
  and ones (\cref{def:adjmatrix}). Its power $A^k$ is the product of $k$
  copies of $A$, with $A^1=A$, and $A^0=I$ is the identity matrix, which has
  $1$ on the diagonal and $0$ elsewhere.
  \item $k$ is a nonnegative integer used for the length of a walk and for
  the exponent of $A$.
  \item $d(n_i)$ is the degree of $n_i$ (\cref{def:degree}).
  \item $u,v,w$ are arbitrary vertices.
  \item For $N\times N$ matrices $X$ and $Y$ the entry $(XY)_{ij}$ equals
  $\sum_{r=1}^{N}X_{ir}Y_{rj}$.
\end{itemize}

\subsection{Conventions}

The photos, the assignment text and the textbook leave several points open.
Each choice we make is listed here and cited by number afterwards.

\begin{convention}[kind of graph]\label{conv:graph}
All graphs are finite, simple and undirected. Simple means no loops and no
repeated edges. The assignment text states this, and the symmetric matrix with
zero diagonal on the board is consistent with it \cite{whiteboard1}. Directed
graphs and graphs with repeated edges are outside the scope of this note; the
adjacency matrix of a directed graph is not symmetric in general, so the
statements below that use symmetry do not carry over unchanged.
\end{convention}

\begin{convention}[the mark above $A$]\label{conv:power}
The board writes a small mark above $A$ in Q1. In the two photos it looks like
the letter $n$ or $k$ set between quotation-like marks, so the board alone does
not settle the reading \cite{whiteboard1,whiteboard2}. The assignment text asks
about $A^k$, so we read the mark as an exponent $k$. A prime or quotation marks
would ask about another matrix or only name $A$.
\end{convention}

\begin{convention}[the board's letters]\label{conv:edgeset}
The board writes $(G,V)$ for the graph and ``$e_{ij}\in V$'' in the rule
that defines $A_{ij}$ \cite{whiteboard1,whiteboard2}. The textbook writes
$G=(V,E)$ and the same rule with $e_{ij}\in E$ \cite[p.~9]{petrosyan2015graph}.
We read the board's $V$ in that rule as the edge set $E$. This is a difference
of notation, and we write $E$.
\end{convention}

\begin{convention}[matrix against drawing]\label{conv:matrix}
The board's drawing and the board's matrix disagree, edge by edge
(\cref{sec:board-mismatch}). The problem statement on the board says the graph
is represented through its adjacency matrix \cite{whiteboard1}, so the matrix is
the example graph. We read the arrowhead on one line of the drawing as a stray
mark, because the matrix is symmetric and the graph is undirected.
\end{convention}

\begin{convention}[the word ``path'']\label{conv:path}
The board does not define path \cite{whiteboard1}. The assignment text says
that in the course a path is a walk with no repeated vertex. The textbook has
the same object under the name \emph{simple path} (Definition~1.2.13), and
it uses the word \emph{path} for a walk with no repeated edge
(Definition~1.2.12) \cite[p.~22]{petrosyan2015graph}. Because the two meanings
of ``path'' differ, we write \emph{trail} for the textbook's Definition~1.2.12
and say ``the textbook's path'' when we quote it. Plain ``path'' always means
no repeated vertex. \Cref{tab:conventions} compares the names.
\end{convention}

\begin{convention}[walks of length zero]\label{conv:zero}
A single vertex is a walk of length $0$ from that vertex to itself. The
textbook does not say whether length $0$ is allowed in its walk definition on
the pages we viewed (printed pp.~21 to 23) \cite{petrosyan2015graph}. Grinberg
allows it \cite[p.~36, Definition~2.9.1(a)]{grinberg2025graph}. We adopt it as a
convention. It makes $A^0=I$ count walks of length $0$, and it makes the
walk-to-path lemma (\cref{lem:walkpath}) true for $u=v$.
\end{convention}

\begin{convention}[page numbers]\label{conv:pages}
Textbook page numbers are the printed numbers; the page of the PDF file is one
higher. The textbook is in Armenian; we give its terms in English translation, and
every definition is cited by the book's number.
\end{convention}

This note makes no claim about what the lecturer intended beyond what the
photos and the assignment text show.
